\ifx\danslelivre\undefined
\documentclass[10pt]{../fichiers-configuration-latex/smfart}
\usepackage{ifthen}
\newboolean{english}
\newboolean{francais}
\setboolean{english}{false}
\setboolean{francais}{true}
\input{../fichiers-configuration-latex/paquets}
\selectlanguage{francais}
\input{../fichiers-configuration-latex/lemme-proposition-theoreme}
\input{../fichiers-configuration-latex/numerotation}
\input{../fichiers-configuration-latex/macros}

\author{Luc Illusie}


\begin{document}

\begin{center}
VI Log régularité, actions très modérées\\
Luc Illusie
\end{center}

\else
\setboolean{english}{false}
\setboolean{francais}{true}
\selectlanguage{francais}
\chapter{Log régularité, actions très modérées}
\addtocontents{toc}{par Luc Illusie}
\begin{center}
Luc Illusie
\end{center}
\fi

\section{Log régularité} \subsection{}\label{6-1.1} Pour le langage des log schémas nous renvoyons
le lecteur à \cite{LSFI@Kato}, \cite{Toric@Niziol}, \cite{HT@GR}. Sauf mention du contraire, les log
structures considérées le seront au sens de la topologie étale. Un log schéma \textit{fin} (resp.
\textit{fs}, i. e. \textit{fin et saturé}) \cite{LSFI@Kato} est un schéma muni d'une log structure
admettant localement (pour la topologie étale) une carte sur un monoïde fin (resp. fin et saturé).
On note en général $M_X$ le faisceau de monoïdes d'un log schéma $X$, $\alpha : M_X \ra
\OO_X$ le morphisme structural, et $\overline{M}_X = M_X/\OO^*_X$. Sauf mention du
contraire, les log schémas considérés sont supposés localement noethériens.

Dans ce qui suit, $X$ désigne un fs log schéma.  \subsection{}\label{6-1.2} Soient $\overline{x}$ un
point géométrique de $X$, d'image $x \in X$, $\OO_{X,\overline{x}}$ le localisé strict de
$X$ en $\overline{x}$. Notons $I_{\overline{x}}$ l'idéal de $\OO_{X,\overline{x}}$ engendré
par $\alpha(M_{X,\overline{x}}-\OO_{X,\overline{x}}^*)$, $C_{X,\overline{x}}$ le sous-schéma
fermé de $X_{(\overline{x})} = \Spec \OO_{X,\overline{x}}$ défini par $I_{\overline{x}}$.
$C_{X,\overline{x}}$ est la trace sur $X_{(\overline{x})}$ de la strate de $X$ où le rang de
$\overline{M}^{\mathrm{gp}}$ est égal à $r(x) =
\mathrm{rg}(\overline{M}_{X,\overline{x}}^{\mathrm{gp}})$. 

On dit que $X$ est \textit{log régulier} en $x$ (ou $\overline{x}$) si $C_{X,\overline{x}}$ est
régulier et de codimension égale à $r(x)$ (cette condition ne dépend que de $x$). On dit que $X$ est
\textit{log régulier} si $X$ est log régulier en tout point. La définition analogue pour les log
schémas zariskiens est due à Kato \cite{Toric@Kato}. La variante dans le cadre étale a été traitée
par Niziol \cite{Toric@Niziol}. Nous rappelons ci-après quelques propriétés de cette notion.
\subsection{}\label{6-1.3} $X$ est log régulier en $x$ si et seulement si $X$, muni de la log
structure Zariski $M_X^{\mathrm{Zar}} := \varepsilon_*M_X$, où $\varepsilon : X_{\mathrm{\et}} \ra
X_{\mathrm{Zar}}$, est log régulier en $x$ au sens de Kato \cite{Toric@Kato},
(\cite[4.6]{Saturated@Tsuji}, \cite[2.4]{Toric@Niziol}. En particulier, si  $X$ est log régulier en
$x$, $X$ est log régulier en toute générisation $y$ de $X$ (\cite[7.1]{Toric@Kato}).

Si la log structure de $X$ est triviale, $X$ est log régulier si et seulement si $X$ est régulier au
sens usuel.  \subsection{}\label{6-1.4} Supposons $X$ log régulier. Soit $j : U \inj X$ l'inclusion
de l'ouvert de trivialité de sa log structure. Alors $U$ est un ouvert dense de $X$ et on a $$ M_X =
\OO_X \cap j_*\OO_U^* $$ (\cite[2.6]{Toric@Niziol}).

Nous dirons qu'un couple $(X,Z)$ formé d'un schéma $X$ et d'un fermé $Z$ est un couple \textit{log
régulier} si, pour la log structure sur $X$ définie par $M_X = \OO_X \cap
j_*\OO_U^*$, où $j : U \inj X$ est l'ouvert complémentaire de $Z$, $X$ est log régulier et
$Z$ est le complément de l'ouvert de trivialité de sa log structure. La log structure précédente sur
$X$ sera dite \textit{associée} au couple $(X,Z)$.  \subsection{}\label{6-1.5} Supposons $X$ log
régulier. Pour $i \in \NN$, soit $X^{(i)}$ l'ensemble des points $x$ de $X$ tels que $r(x) =
i$, avec la notation de \ref{6-1.2}. C'est une partie localement fermée, sous-jacente à un
sous-schéma régulier de $X$, de codimension $i$, dont la trace sur $X_{(\overline{x})}$, en chaque
point géométrique $\overline{x}$ localisé en  $x \in X^{(i)}$, est $C_{X,\overline{x}}$. On dit que
$X^{(i)}$ est la \textit{strate de codimension} $i$ définie par le rang de
$\overline{M}^{\mathrm{gp}}$. La stratification par les $X^{(i)}$ est appelée \textit{stratification
par le rang de $\overline{M}^{gp}$}, ou \textit{stratification canonique}. Voici deux exemples. 

(i) Si $X$ est un schéma noethérien régulier, muni de la log structure définie par un diviseur à
croisements normaux $D$, $X^{(i)}$ est l'ensemble des points où passent exactement $i$ branches de
$D$, i. e. tels que le normalisé de $D$ ait $i$ points au-dessus de $x$.

(ii) Si $X$ est une variété torique sur un corps $k$, de tore $T$, munie de sa log structure
canonique, $X$ est un log schéma log régulier, l'ouvert de trivialité de la log structure est $T$,
et $X^{(i)}$ est la réunion des orbites de $T$ de codimension $i$.  \subsection{}\label{6-str loc
log rég}\label{6-str loc log reg} Supposons $X$ log régulier en $x$, image du point géométrique
$\overline{x}$, soient $P = \overline{M}_{X,\overline{x}}$, $k = k(\overline{x)}$. Notons que $P$
est un monoïde fs \textit{saillant} (i. e. tel que $P^* = 0$). Soit $\widehat  X_{\overline{x}}$ le
complété de $X_{(\overline{x})}$ au point fermé. Alors, d'après \cite[3.2]{Toric@Kato}, $X$ admet
une carte modelée sur $P$ en $\overline{x}$, qui donne lieu à des isomorphismes

(i) $$ \widehat  X_{\overline{x}} \iso \Spec k[[P]][[t_1,\cdots,t_n]] $$ si $\OO_{X,x}$ est
d'égale caractéristique,

(ii) $$ \widehat  X_{\overline{x}} \iso \Spec C(k)[[P]][[t_1,\cdots,t_n]]/(f) $$ si
$\OO_{X,x}$ est d'inégale caractéristique $(0,p)$, où $C(k)$ est un anneau de Cohen pour k,
et $f$ est congru à $p$ modulo l'idéal engendré par $P- \{0\}$ et les $t_i$.
\subsection{}\label{6-1.7} Supposons $X$ log régulier. Alors $X$ est régulier en $x$, image de
$\overline{x}$, si et seulement si $\overline{M}_{\overline{x}} \isononcan \NN^r$
(\cite[5.2]{Toric@Niziol}, voir aussi \cite [1.8]{LR@Vidal}). Il en résulte que l'ensemble des
points de régularité de $X$ coïncide avec l'ensemble des points de régularité du monoïde
$\overline{M}_{\overline{x}}$, et en particulier est ouvert dans $X$ (\cite[5.3]{Toric@Niziol}). Si
$X$ est log régulier et régulier, l'ouvert de trivialité de la log structure est alors le complément
d'un diviseur à croisements normaux.  \subsection{} Soit $f : X \ra Y$ un morphisme de log schémas
fs. Si $Y$ est log régulier et $f$ log lisse, $X$ est log régulier.

L'analogue pour les log structures Zariski est \cite[8.2]{Toric@Kato}. La démonstration de
(\textit{loc. cit.}) s'applique, \textit{mutatis mutandis}, dans le cadre étale.

Le corollaire suivant jouera un rôle clé dans l'application des résultats de de Jong à la
démonstration du théorème d'uniformisation de Gabber. Rappelons que, si $S$ est un schéma, une
\textit{courbe nodale} $f : C \ra S$ est un morphisme propre et plat, purement de dimension relative
1, dont les fibres géométriques ont pour seules singularités des croisements normaux.
\begin{proposition} \label{6-nodal implique log lisse} Soit $(Y,T)$ un couple log régulier
(\ref{6-1.4}). Soit $f : X \ra Y$ une courbe nodale, lisse au-dessus de $Y-T$. Soit $D$ un diviseur
effectif sur $X$, de support contenu dans le lieu de lissité de $f$, et étale sur $Y$. Alors le
couple $(X,f^{-1}(T) \cup D)$ est log régulier, et pour les log structures associées, $f$ est un
morphisme de log schémas et est log lisse.  \end{proposition} La question est locale sur $X$ et
l'assertion est triviale sur l'ouvert de lissité de $f$ et sur $f^{-1}(Y-T)$. Soit $\overline{x}$ un
point géométrique de non lissité de $f$, d'image $\overline{y}$ dans $Y$. D'après \sga{7}{XV}{1.3.2}
(voir aussi \cite[2.23]{Alterations@de_Jong}, on a \begin{equation} \OO_{X,\overline{x}}
\simeq \OO_{Y,\overline{y}}\{u,v\}/(uv-h), \label{6-(1.9.1)} \end{equation} où
$\OO_{Y,\overline{y}}\{u,v\}$ désigne le localisé strict de $\Spec
\OO_{Y,\overline{y}}[u,v]$ au point $(u = v = 0)$ au-dessus du point fermé de
$Y_{(\overline{y})}$, et $h$ est un élément de $\OO_{Y,\overline{y}}$ inversible sur $Y-T$,
donc appartenant à $M_{Y,\overline{y}}$ (\ref{6-1.4}). Soit $c$ une carte locale de $Y$ en
$\overline{y}$ définie par une section de  $M_{Y,\overline{y}} \ra P :=
\overline{M}_{Y,\overline{y}}$. Écrivons $h = \varepsilon a$, avec $\varepsilon \in
\OO^*_{Y,\overline{y}}$ et $a \in P$. Soit $Q$ le monoïde fs défini par le carré cocartésien
$$ \xymatrix{\NN^{2} \ar[r] & Q \\ \NN \ar[u] \ar[r] & \ZZ \times P
\ar[u]^{g}}, $$ où la flèche $\NN \ra \NN^2$ (resp. $\NN \ra \ZZ \times
P $) est donnée par $1 \mapsto (1,1)$ (resp. $1 \mapsto (1,a)$). Le carré $$ \xymatrix{\Spec
\OO_{Y,\overline{y}}[u,v]/(uv-h) \ar[r]^-{d} \ar[d] & \Spec \ZZ[Q] \ar[d]^{\Spec
\ZZ[g]} \\ \Spec \OO_{Y,\overline{y}} \ar[r]^{c'} & \Spec \ZZ[\ZZ\times
P]}, $$ où la flèche $c'$ est donnée par $c$ et $\ZZ \ra \OO_{Y,\overline{y}}$, $1
\mapsto \varepsilon$, et $d$ par $(1,0) \mapsto u$, $(0,1) \mapsto v$ sur $\NN^2$ et $c'$ sur
$\ZZ \times P$, est cartésien. Comme $c'$ est une carte de $Y$ en $\overline{y}$, $d$ est une
carte en $\overline{x}$ de $X$ muni de la log structure image inverse par $f$, i. e. définie par
$\OO_X \cap \OO^*_{f^{-1}(Y-T)}$, et $(c',d,g)$ est une carte de $f$. Comme $g$ est
injectif et $\coker g^{gp} \simeq \ZZ$, $f$ est log lisse en $\overline{x}$, et le couple
$(X,f^{-1}(T) \cup D)$ est log régulier en $\overline{x}$.  \section{Revêtements Kummer étales} Dans
cette section et la suivante, nous aurons à considérer des action de groupes sur des schémas ou des
log schémas : sauf mention du contraire, il s'agira d'actions \textit{à droite}.
\subsection{}\label{6-Kummer} Rappelons quelques définitions (cf. \cite[3]{Overview@Illusie},
\cite[3.1]{Ket@Stix}, \cite{GFL@Vidal}). Un homomorphisme $h : P \ra Q$ de monoïdes intègres  est
dit \textit{de Kummer} si $h$ est injectif et, pour tout $q \in Q$, il existe $n \ge 1$ et $p \in P$
tel que $nq = h(p)$. On dit qu'un morphisme $f : X \ra Y$ de fs log schémas est \textit{de Kummer}
si, pour tout point géométrique $\overline{x}$ de $X$ d'image $\overline{y} = f(\overline{x})$ dans
$Y$, l'homomorphisme induit $\overline{M}_{\overline{y}} \ra \overline{M}_{\overline{x}}$ est de
Kummer. On dit qu'un morphisme $f : X \ra Y$ est \textit{Kummer étale} si $f$ est de Kummer et log
étale. On dit que $f$ est un \textit{revêtement Kummer étale} si $f$ est de Kummer étale et le
morphisme de schémas sous-jacent est fini.

Le \textit{site Kummer étale} d'un fs log schéma $X$ est la catégorie des fs log schémas de Kummer
étales au-dessus de $X$ munie de la topologie définie par les familles surjectives de $
X$-morphismes (lesquels sont automatiquement de Kummer étale). Pour $X$ connexe,  les revêtements
Kummer étales de $X$ forment une catégorie galoisienne, équivalente à la catégorie des
représentations d'un groupe profini, le \textit{groupe fondamental logarithmique} de $X$,
$\pi_1^{\mathrm log}(X,x)$, où $x$ est un \textit{point géométrique logarithmique} de $X$, cf.
(\textit{loc. cit.}).  \subsection{}\label{6-ket standard} Un morphisme $f : X \ra Y$ de log schémas
fs qui est déduit par changement de base par un morphisme strict $g : Y \ra \Spec \ZZ[P]$ d'un
morphisme $\Spec \ZZ[h] :\Spec \ZZ[Q] \ra \Spec \ZZ[P]$, où $h : P \ra Q$ est un homomorphisme de Kummer
entre monoïdes fs tels que le conoyau de $h^{\mathrm{gp}}$ soit annulé par un entier $n$ inversible
sur $Y$, est un revêtement Kummer étale. On dit qu'un tel revêtement est un revêtement Kummer étale
\textit{standard}. 

Un revêtement Kummer étale standard $f : X \ra Y$ comme ci-dessus est \textit{galoisien} dans le
sens suivant : le groupe diagonalisable (étale) $G =
Hom(Q^{\mathrm{gp}}/P^{\mathrm{gp}},\mathbf{G}_{m}{}_{Y})$ opère sur $X$ par automorphismes de
$Y$-log schémas, et le morphisme canonique $$ X \times_Y G \ra X \times_Y X, (x,g) \mapsto (x,xg),
$$ où le produit au second membre est pris dans la catégorie des fs log schémas, est un isomorphisme
(cf. \cite[3.2]{Overview@Illusie}). 

Tout morphisme de Kummer étale est, localement pour la topologie étale, isomorphe à un revêtement
Kummer étale standard (cf. \cite[3.1.4]{Ket@Stix}). Plus précisément, si $f : X \ra Y$ est un
morphisme de Kummer étale surjectif, avec $X$ (resp. $Y$) strictement local de point fermé $x$
(resp. $y$), on a un carré cartésien $$ \xymatrix{X \ar[r]^-a \ar[d]^f & \Spec \ZZ[Q] \ar[d]^{\Spec
\ZZ[h]} \\ Y \ar[r]^-b & \Spec \ZZ[P]}, $$ où $a$ (resp. $b$) est une carte de $X$ (resp. $Y$), et $h
: P \ra Q$ un morphisme de Kummer tel que $\coker h^{\mathrm{gp}}$ soit annulé par un entier $n$
inversible sur $Y$. Alors $f$ est un revêtement Kummer étale galoisien de groupe $G =
\mathrm{Hom}(\coker h^{gp},\mu_n(k(y)))$.

\subsection{}\label{6-Kummer standard local} Dans ce cas, l'action de $G$ sur $X$ peut se décrire de
la manière suivante. Soit $C = \coker h^{gp}$. La suite exacte $$ 0 \ra P^{\mathrm{gp}} \ra
Q^{\mathrm{gp}} \ra C \ra 0 $$ donne, par application du dual de Cartier
$\mathrm{Hom}(-,\GG_m)$, une suite exacte de groupes diagonalisables $$ 0 \ra \Gamma \ra T_Q
\ra T_P \ra 0, $$ où $T_P$ (resp. $T_Q$) est le tore sur $\ZZ$ dual de $P$ (resp. $Q$), et $G$ est
la valeur du groupe fini constant $\Gamma_Y$. Posons $Z_P = \Spec \ZZ[P]$, $Z_Q = \Spec \ZZ[Q]$.
L'accouplement canonique $$ T_P \otimes P^{\mathrm{gp}} \ra \GG_m $$ définit une famille de
caractères $$ (\chi_p : T_P \ra \GG_m)_{p \in P^{\mathrm{gp}}}, $$ qui détermine l'action de
$T_P$ sur $Z_P$ par $$ g.p = \chi_p(g)p, $$ pour un point $g$ de $T_P$ (à valeurs dans $S$), et $p
\in P$, vu comme un point de $Z_P$ à valeurs dans $S$. L'action de $T_Q$ sur $Z_Q$ est décrite de
façon similaire, et la carte $a : X \ra Z_Q$ est équivariante relativement aux actions de $G$ et
$T_Q$ sur $X$ et $Z_Q$ respectivement : pour $g \in G$ et $q \in Q$, \begin{equation} g.a^*(q) =
\chi_q(g)a^*(q), \label{6-g.q} \end{equation} où $a^*(q) \in M_{X,x}$ est l'image de $q$ par $a$. En
particulier, $G$ opère \textit{trivialement} sur $\overline{M}_{X,x}= M_{X,x}/\OO^*_{X,x}$. 

Par ailleurs, $P = Q \cap P^{\mathrm{gp}}$, donc $$ P = \{q \in Q ; \chi_q(g) = 1 \  \forall g \in G
\}.  $$ Ainsi, $Y$ est un quotient de $X$ par $G$, en tant que log schéma, et en tant que schéma :
$$ \OO_Y = (f_*\OO_X)^G \ \ ; \ \ M_Y = (f_*M_X)^G.  $$

Supposons de plus $Y$ log régulier, de sorte que $X$ l'est également. Alors $G$ opère librement
au-dessus de l'ouvert (dense) de $Y$ où la log structure est triviale, et fait de $X$ un revêtement
étale galoisien de $Y$ de groupe $G$. De plus, $G$ opère \textit{trivialement} sur la
\textit{strate} $C_{X,x}$ (\ref{6-1.5}). Pour le voir, on peut remplacer $X$ (resp. $Y$) par son
complété en $x$ (resp. $y$). On peut supposer de plus $P$ saillant, i. e.  $P^* = \{0\}$ (et même $P
= \overline{M}_{Y,y}$). Alors $Q$ est également saillant : on a $Q = Q^* \oplus Q_1$ avec $Q_1$
saillant, et $Q^* \cap P = \{0\}$, donc comme $Y \times_{\Spec \ZZ [P]} \Spec \ZZ [Q_1]$ est local, $Q^*
= \{0\}$.  D'après \ref{6-str loc log rég}, la carte $(h,a,b)$ de $f$ permet d'identifier
$\OO_{Y,y} \ra \OO_{X,x}$ à un homomorphisme de la forme $$ k[[t_1,\cdots,t_r]][[P]]
\ra k[[t_1,\cdots,t_r]][[Q]], $$ dans le cas d'égale caractéristique (resp.  $$
C[[t_1,\cdots,t_r]][[P]]/(u) \ra C[[t_1,\cdots,t_r]][[Q]]/(u) $$ dans le cas d'inégale
caractéristique $(p,0)$), où $k = k(x)  = k(y)$ (resp. $C$ est un anneau de Cohen de $k$, et $u$ est
congru à $p$ modulo l'idéal engendré par $P - \{0\}$ et les $t_i$), les $t_i$ et $u$ étant fixes par
$G$. Les strates $C_{X,x}$  et $C_{Y,y}$ sont respectivement $\Spec B$ et $\Spec A$, où $A$ et $B$ sont
les réductions de $\OO_{Y,y}$ et $\OO_{X,x}$ modulo les idéaux engendrés par
$P-\{0\}$ et $Q - \{0\}$. Comme $G$ opère par \ref{6-g.q} sur $Q$, $G$ opère donc trivialement sur
$C_{X,x}$, et $f : X \ra Y$ induit un isomorphisme $C_{X,x} \iso C_{Y,y}$.  \section{Actions très
modérées} \subsection{}\label{6-3.1} Soit $X$ un log schéma fs, muni d'une action d'un groupe fini
$G$. On se propose de dégager des conditions suffisantes sur l'action de $G$ pour que, lorsque $X$
est log régulier, le quotient de $X$ par $G$ existe comme log schéma et soit log régulier. 

On dit que $G$ opère \textit{modérément} sur $X$ en un point géométrique $\overline{x}$ de $X$
localisé en $x$, si le stabilisateur $G_{\overline{x}}$ de $\overline{x}$ (\textit{groupe d'inertie
en $x$}) est d'ordre premier à la caractéristique de $k(x)$. On dit que $G$ opère modérément sur $X$
si $G$ opère modérément en $\overline{x}$ pour tout $\overline{x}$. Ces définitions ne font pas
intervenir la log structure de $X$.

La définition et les résultats qui suivent sont dues à Gabber. On dit que $G$ opère \textit{très
modérément} sur $X$ \textit{en} $\overline{x}$ si les trois conditions suivantes sont satisfaites :

(i) $G$ opère modérément en $\overline{x}$ ;

(ii) $G$ opère trivialement sur $\overline{M}_{X,\overline{x}}$ ;

(iii) $G_{\overline{x}}$ opère trivialement sur la strate $C_{X,\overline{x}}$ (\ref{6-1.2}).

On dit que $G$ opère \textit{très modérément} sur $X$ si $G$ opère très modérément sur $X$ en tout
point géométrique.

Le résultat principal est le suivant : \begin{theoreme}\label{6-très modéré}\label{6-tres modere}\label{6-3.2}
Soit $X$ un log schéma fs log régulier, muni d'une action d'un groupe fini $G$. On suppose que $G$
opère de façon admissible sur le schéma sous-jacent à $X$, librement sur un ouvert dense, et très
modérément. Alors le quotient $Y = X/G$ existe comme log schéma, est un log schéma fs log régulier,
et la projection $f : X \ra Y$ est un revêtement Kummer étale de groupe $G$.  \end{theoreme}
Rappelons (\sga{1}{V}{1}) que dire que $G$ opère de façon admissible signifie que $X$ est réunion
d'ouverts affines stables par $G$, de sorte que le quotient $X/G$ est un schéma. 

Compte tenu de \ref{6-ket standard}, on en déduit : \begin{corollaire}\label{6-inertie abélienne}
Sous les hypothèses de \ref{6-très modéré}, pour tout point géométrique $\overline{x}$ de $X$
localisé en $x$, d'image $\overline{y}$ dans $Y$, le groupe d'inertie $H = G_{\overline{x}}$ en $x$
est abélien, d'ordre premier à la caractéristique de $k(x)$. Le localisé strict $X_{(\overline{x)}}$
est un revêtement Kummer étale du localisé strict $Y_{(\overline{y})}$ de groupe $H$, et le
revêtement Kummer étale $Y_{(\overline{y})} \times_Y X$ induit sur $Y_{(\overline{y})}$ s'en déduit
par extension du groupe $H$ à $G$.  \end{corollaire} \subsection{}\label{6-carte équivariante} Nous
démontrerons un résultat plus précis que \ref{6-très modéré}, pour lequel nous aurons besoin de la
notion suivante. Soient $n$ un entier $\ge 1$, $G$ un groupe fini, et $Q$ un monoïde fs. Supposons
donné un homomorphisme $$ \chi : G^{ab} \otimes Q^{gp} \ra \mu_n := \mu_n(\CC), \ \ g \otimes q
\mapsto \chi_q(g).  $$ Soit $\Lambda = \ZZ[\mu_n][1/n]$. On déduit de $\chi$  une action de $G$ sur
le log schéma $\Spec \Lambda[Q]$, caractérisée par $$ g.q = \chi_q(g)q $$ pour $g \in G$, $q \in Q$.

Soit $X$ un log schéma fs muni d'une action de $G$. Par une \textit{carte $G$-équivariante de $X$
modelée sur $\chi$ et Q}, on entend un morphisme strict, $G$-équivariant $$ c : X \ra \Spec
\Lambda[Q].  $$ \begin{proposition}\label{6-très modéré ponctuel}\label{6-tres modere ponctuel} Soit
$X$ un fs-log schéma, muni d'une action d'un groupe fini $G$, et soit $\overline{x}$ un point
géométrique de $X$ localisé en $x$. On note $H = G_{\overline{x}}$ le groupe d'inertie en $x$. 

(a) On suppose que les conditions (i) et (ii) de \ref{6-3.1} sont vérifiées en $\overline{x}$. Soit
$n$ l'exposant de $|H|$. Il existe, localement pour la topologie étale au voisinage de
$\overline{x}$, une carte $H$-équivariante de $X$ modelée sur $Q =\overline{M}_{\overline{x}}$ et un
homomorphisme $\chi : H^{ab} \otimes Q^{gp} \ra \mu_n$. 

(b) On suppose de plus que $G$ agit très modérément en $\overline{x}$, i. e. que $H$ agit
trivialement sur la strate $C_{\overline{x}}(X)$. Alors, localement pour la topologie étale au
voisinage de $\overline{x}$ et de son image $\overline{y}$ dans $Y$, le quotient $Y = X/H$ existe
comme log schéma, et $X$ est un $Y$-sous log schéma fermé strict d'un revêtement Kummer étale
standard de $Y$.

(c) Si, sous les hypothèses de (ii), $X$ est supposé en outre log régulier, alors
$Y_{(\overline{y})}$ est log régulier, et $X_{(\overline{x})}$ est un revêtement Kummer étale
standard de $Y_{(\overline{y})}$ de groupe $H$.  \end{proposition} Sous les hypothèses de
\ref{6-très modéré}, le log schéma de schéma sous-jacent $Y = X/G$, muni de la log structure $M_Y =
(f_*M_{X})^{G}$ est log régulier, et s'identifie, par localisation stricte en $\overline{y}$ au
quotient $Y_{(\overline{y})} = X_{(\overline{x})}/H$. Le fait que $G$ opère génériquement librement
implique la dernière assertion de \ref{6-inertie abélienne}, et \ref{6-très modéré} s'en déduit.

Par ailleurs, compte tenu de la description locale \ref{6-Kummer standard local} des revêtements
Kummer étales standard, \ref{6-très modéré ponctuel} implique : \begin{corollaire} \label{6-tres
modere dans voisinage} Sous les hypothèses de \ref{6-très modéré ponctuel}, si $X$ est log régulier
et si $G$ agit très modérément en $\overline{x}$, alors $G$ agit très modérément sur $X$ dans un
voisinage étale de $\overline{x}$.  \end{corollaire} \subsection{} \label{6-preuve très modéré
ponctuel} \label{6-preuve tres modere ponctuel} \textit{Preuve de  \ref{6-très modéré ponctuel}} 

(a) Considérons la suite exacte de groupes abéliens \begin{equation} 0 \ra
\OO_{X,\overline{x}}^* \ra M_{X,\overline{x}}^{gp} \ra \overline{M}_{X,\overline{x}}^{gp}
\ra 0.  \label{6-exacte} \end{equation} Elle est $H$-équivariante, et si $Q =
\overline{M}_{\overline{x}}$, $Q^{gp}$ est de type fini sans torsion. Choisissons un scindage $s :
Q^{gp} \ra M_{X,\overline{x}}^{gp}$ de \ref{6-exacte} (comme suite exacte de groupes abéliens).
Comme $H$ agit trivialement sur $Q^{gp}$, on a, pour $a \in Q^{gp}$ et $g \in H$ $$ g.s(a) =
z(g,a)s(a) $$ avec $z(g,a) \in \OO_{X,\overline{x}}^*$. Pour $g_1$, $g_2$ dans $H$, on a $$
z(g_1g_2,a) = (g_1z(g_2,a)).z(g_2,a), $$ en d'autres termes, $g \mapsto (a \mapsto z(g,a))$ est un
1-cocycle $$ z \in Z^1(H,\mathrm{Hom}(Q^{gp},\OO_{X,\overline{x}}^*)).  $$ L'image $[z]$ de
$z$ dans $H^1(H,\mathrm{Hom}(Q^{gp},\OO_{X,\overline{x}}^*))$ est la classe de cohomologie
de \ref{6-exacte}. Dans le carré commutatif de flèches canoniques $$
\xymatrix{Z^1(H,\mathrm{Hom}(Q^{gp},\OO_{X,\overline{x}}^*)) \ar[r] \ar[d]
&Z^1(H,\mathrm{Hom}(Q^{gp},k(\overline{x})^*)) \ar[d] \\
H^1(H,\mathrm{Hom}(Q^{gp},\OO_{X,\overline{x}}^*)) \ar[r] &
H^1(H,\mathrm{Hom}(Q^{gp},k(\overline{x})^*))}, $$ la flèche verticale de droite est (trivialement)
un isomorphisme. Comme $n$ est premier à la caractéristique de $k(x)$,  $z$ est l'unique relèvement
de son image $\chi$ dans $Z^1(H,\mathrm{Hom}(Q^{gp},k(\overline{x})^*))$, et la flèche horizontale
inférieure est un isomorphisme, puisque $(1 + \mathbf{m}_{X,\overline{x}})^*$ est $n$-divisible.
Identifiant $Z^1(H,\mathrm{Hom}(Q^{gp},k(\overline{x})^*))$ à
$\mathrm{Hom}(H,\mathrm{Hom}(Q^{gp},k(\overline{x})^*))$, on trouve ainsi un homomorphisme $$ \chi :
H^{ab} \otimes Q^{gp} \ra \mu_n, $$ et la restriction de $s$ à $Q$ est, au voisinage de
$\overline{x}$, une carte $H$-équivariante $a$ de $X$ modelée sur $Q$ et $\chi$.

(b) On peut remplacer $X$ par son localisé strict $X_{(\overline{x})}$ et $G$ par $H$. Le quotient
$Y = X/G$ existe alors comme schéma, est strictement local, de point fermé $\overline{y}$, avec
$k(\overline{y}) = k(\overline{x})$, et $X \ra Y$ est fini. Soit $P'$ le sous-groupe de $Q^{gp}$
défini par $$ P' = \{q \in Q^{gp} | \chi_q(g) = 1 \ \ \forall g \in G\}.  $$ Il est d'indice fini,
premier à la caractéristique de $k(\overline{x})$. Soit $$ P = P' \cap Q.  $$ La carte équivariante
$a : X \ra \Spec \Lambda[Q]$ de (a) (où $\Lambda = \ZZ [\mu_n][1/n]$) définit, par passage au
quotient, un morphisme $b : Y \ra \Spec \Lambda[P]$. Ce morphisme $b$ définit une log structure sur
$Y$ pour laquelle $(f_*M_X)^G = M_Y$ ($f : X \ra Y$ désignant la projection), et une carte de $f$
\begin{equation} \xymatrix{X \ar[r] \ar[d]^f & \Spec \Lambda[Q] \ar[d] \\ Y \ar[r] & \Spec \Lambda[P]},
\label{6-carré} \end{equation} modelée sur l'homomorphisme de Kummer $P \subset Q$. Comme $G$ agit
trivialement sur la strate $C_{\overline{x}}(X) = \Spec
\OO_{X,\overline{x}}/I_{\overline{x}}$, et que $G$ est d'ordre premier à la caractéristique
de $k(\overline{x})$, l'homomorphisme $$ \OO_{Y,\overline{y}} (=
\OO_{X,\overline{x}}^G) \ra \OO_{X,\overline{x}}/I_{\overline{x}} (=
(\OO_{X,\overline{x}}/I_{\overline{x}})^G) $$ est surjectif, donc il en est de même de
l'homomorphisme $$ k(\overline{y})[Q] \  = \ k(\overline{x})[Q] \ra
\OO_{X,\overline{x}}/\mathbf{m}_{Y,\overline{y}}\OO_{X,\overline{x}}.  $$ Par
Nakayama, il en résulte que le morphisme \begin{equation} i : X \ra Y \times_{\Spec \Lambda[P]} \Spec
\Lambda[Q] \label{6-morphisme} \end{equation} déduit de \ref{6-carré} est une immersion fermée
stricte.

(c) Dans ce cas, la strate $C_{\overline{x}}(X)$ est régulière, et se projette isomorphiquement sur
$C_{\overline{y}}(Y)$. Comme $\mathrm{rg}(\overline{M}^{\mathrm{gp}}_{\overline{x}}) =
\mathrm{rg}(\overline{M}^{\mathrm{gp}}_{\overline{y}}) =
\mathrm{codim}(C_{\overline{y}}(Y),Y_{(\overline{y})})$, $Y$ est log régulier en $y$, et l'immersion
$i$ est un isomorphisme, de sorte que $X_{(\overline{x})}$ est un revêtement Kummer étale de
$Y_{(\overline{y})}$ de groupe $H$.  \section{Points fixes} \subsection{}\label{6-stratification}
Soit $X$ un schéma noethérien séparé sur lequel opère un groupe fini $G$. Pour chaque sous-groupe
$H$ de $G$, on note $X^H$ le schéma des points fixes de $H$. C'est un sous-schéma fermé de $X$,
représentant le foncteur $S \mapsto X(S)^H$, intersection des graphes des translations  $h : X \ra
X$ pour $h \in H$. En chaque point géométrique $\overline{x}$ de $X^H$ d'image $x$, le sous-groupe
d'inertie $G_{\overline{x}}$ de $G$ contient $H$, et est égal à $H$ si et seulement si $x$
appartient au sous-schéma (localement fermé) $$ X_H = X^H - \cup_{H' \supset H, H' \ne H}X^{H'}.  $$
Pour $g \in G$, on a $$ gX^H = X^{gHg^{-1}}, $$ et de même $$ gX_h = X_{gHg^{-1}} $$ de sorte que la
réunion $X^C$ (resp. $X_C$) des $X^H$ (resp. $X_H$) pour $H$ dans une classe de conjugaison $C$ de
sous-groupes de $G$ est $G$-stable. Les $X_C$, pour $C$ parcourant les classes de conjugaison de
sous-groupes de $G$, forment une stratification de $X$ par des sous-$G$-schémas, avec la propriété
que pour tout point géométrique $\overline{x}$ localisé en $X_C$, le groupe d'inertie
$G_{\overline{x}}$ appartient à $C$. Nous appellerons cette stratification la \textit{stratification
par l'inertie}.

Le but de ce numéro est de donner des exemples d'actions très modérées de groupes finis $G$ sur des
log schéma log réguliers et réguliers $X$, où un raffinement de la stratification par l'inertie est
déduite de la stratification canonique d'un diviseur à croisements normaux $G$-stable.

Le résultat suivant est classique, nous en donnons une démonstration, faute de référence.
\begin{proposition}\label{6-pts fixes reguliers} Soit $X$ un schéma noethérien régulier sur lequel
opère un groupe fini $G$ de façon modérée \ref{6-3.1}. Alors le schéma des points fixes $X^G$ est
régulier.  \end{proposition} On peut supposer $X$ local, $X = \Spec A$, de sorte que $X^G = \Spec A_G$,
où $A_G$ est l'algèbre des co-invariants $A/I$, $I$ l'idéal engendré par les $ga - a$, pour $g$ dans
$G$ et $a$ dans l'idéal maximal $\mathbf{m}$ de $A$. On peut supposer de plus que $G$ est contenu
dans le groupe d'inertie au point fermé, sinon $X^G$ est vide. Comme $\widehat X^G = \Spec \hat A/I
\hat A = (\hat X)^G$, on peut supposer $X $local et complet. On démontre le lemme en
\textit{linéarisant} l'action de $G$.

Supposons d'abord que $G$ soit d'égale caractéristique $p$. Choisissons une base $t = (t_i)_{1 \le i
\le r}$ de l'espace cotangent $T = \mathbf{m}/\mathbf{m}^2$ et des éléments $x_i$ dans $\mathbf{m}$
relevant les $t_i$. Choisissons un corps de représentants, noté encore $k$, de $k$ dans $A$, et
notons $$ \varphi : k[[T]] \ra A $$ l'homomorphisme envoyant $t_i$ sur $x_i$. L'homomorphisme $$ f :
k[[T]] \ra A $$ envoyant $t$ sur le système $$ y = (1/\mathrm{card}(G))\sum_{g \in G} g\varphi
g^{-1}t $$ est $G$-équivariant. En outre, $y$ est congru à $x$ modulo $\mathbf{m}^2$, donc est un
système régulier de paramètres de $A$, et donc $f$ est un isomorphisme. Par suite $f$ induit un
isomorphisme $$ k[[T]]_G = k[[T_G]] \ra A_G, $$ où $T_G$ est l'espace des co-invariants de $G$ sur
$T$, ce qui démontre le lemme dans ce cas.

Supposons maintenant que $A$ soit de caractéristique mixte $(0,p)$. Soit $C$ un anneau de Cohen pour
$k$. Choisissons un homomorphisme $C \ra A$ relevant $C \ra k$. Comme $G$ est d'ordre premier à $p$,
il existe un unique $C[G]$-module $V$, libre de type fini sur $C$, relevant $T$. Soit $t = (t_i)_{1
\le i \le r}$ une base de $T$, $v = (v_i)$ une base de $V$ relevant $t$. Comme plus haut,
choisissons des relèvements $x_i$ dans $\mathbf{m}$ des $t_i$, qui forment donc une suite régulière
de paramètres dans $A$. Prolongeons l'homomorphisme $C \ra A$ en $$ \varphi : C[[V]] \ra A $$ en
envoyant $v_i$ sur $x_i$. L'homomorphisme $$ f : C[[V]] \ra A $$ envoyant $v = (v_i)$ sur $$ y =
(1/\mathrm{card}(G))\sum_{g \in G} g\varphi g^{-1}v $$ est $G$-équivariant, et $y$ est congru à $x$
modulo $\mathbf{m^2}$, donc est un système régulier de paramètres de $A$. L'homomorphisme $f$ est
donc surjectif, et son noyau est défini par un élément $F$ congru à $p$ modulo l'idéal engendré par
les $v_i$. Ainsi $X$ est un diviseur régulier $G$-équivariant dans $X' = \Spec C[[V]]$, d'équation
$F=0$. L'image de $F$ dans $C[[V_G]]$ est un paramètre régulier, donc, comme $X'^G = \Spec C[[V_G]]$,
$X^G = X \times_{X'} X'^G$ est régulier.  \begin{corollaire}\label{6-stratification régulière} Soit
$X$ un schéma noethérien séparé, régulier, muni d'une action modérée d'un groupe fini $G$. Alors la
stratification de $X $par l'inertie est formée de schémas réguliers.  \end{corollaire}
\subsection{}\label{6-dcn stables} Sous les hypothèses de \ref{6-stratification régulière}, soit $Y$
un diviseur à croisements normaux stricts $G$-stable, réunion de composantes irréductibles $Y_i$,
$1\le i \le m$. On munit $X$ de la log structure définie par $Y$. Rappelons (cf.
\cite[7.1]{Alterations@de_Jong}) qu'on dit que $Y$ est \textit{$G$-strict} si la condition suivante
est réalisée :  pour tout $i$ et pour tout $g \in G$, si $gY_i \cap Y_i \ne \emptyset$, alors $gY_i
= Y_i$. Si $Y$ est $G$-strict, alors la condition (ii) de \ref{6-3.1} est vérifiée en chaque point
géométrique $\overline{x}$ de $X$. En effet, si $(D_i)_{1 \le i \le r}$ est l'ensemble des branches
de $Y$ passant par $\overline{x}$, alors $gD_i = D_i$ pour tout $i$. Comme
$\overline{M}^{\mathrm{gp}}_{\overline{x}} = \oplus_{1 \le i \le r}\ZZ e_i$, $e_i$ correspondant à
$D_i$, le groupe d'inertie $G_{\overline{x}}$ opère trivialement sur
$\overline{M}^{\mathrm{gp}}_{\overline{x}}$. Rappelons également (\cite[7.2]{Alterations@de_Jong})
qu'il existe une modification $G$-équivariante canonique $f : \tilde X \ra X$ telle que
$f^{-1}(Y)_{\mathrm{red}}$ soit $G$-strict.

\begin{corollaire}\label{6-dcn stricts} Supposons que $Y$ soit $G$-strict, que $G$ opère de façon
modérée, admissible et génériquement libre, et que la stratification canonique de $X$ soit plus fine
que la stratification par l'inertie. Alors $G$ opère très modérément sur $X$ (et donc la conclusion
de \ref{6-très modéré} s'applique). Le groupe d'inertie le long de la strate $X^{(i)}$ (\ref{6-1.5}
(i)) est constant, de valeur $G_i$, avec $\mathrm{rg}(G_i) = \mathrm{codim}(X^{(i)},X)$. En
particulier, $G$ opère librement sur la strate $X^{(0)} = X-Y$.  \end{corollaire}
\subsection{}\label{6-exemples}\textit{Exemples}. 

(a) Soient $k$ un corps algébriquement clos, $n$ un entier $\ge 2$ premier à la caractéristique
résiduelle de $k$, $G$ le groupe $\mu_n = \mu_n(k)$. On fait opérer $G$ sur $X = \mathbf{A}^2_k$ par
homothéties ($(\lambda,x) \mapsto \lambda x$ pour $\lambda \in G$, $x \in X(k)$). La stratification
par l'inertie comporte deux strates, $X-\{0\}$, où $G$ opère librement, et $X^G = \{0\}$. La donnée
de deux droites $Y_1$, $Y_2$ telles que $Y_1 \cap Y_2 = \{0\}$ définit un diviseur à croisements
normaux $G$-strict $Y = Y_1 \cup Y_2$, et le couple $(X,Y)$ vérifie les conditions de \ref{6-dcn
stricts}. Le choix de paramètres $t_1$, $t_2$ tels que $Y_i = V(t_i)$ permet de définir une carte
équivariante \ref{6-carte équivariante} $c : \Spec \ZZ[\NN^2] \la X$, $e_i \mapsto t_i$, associée à
l'homomorphisme $\chi : G \otimes \ZZ^2 \ra \mu_n$ tel que $\chi(\lambda \otimes e_i) = \lambda$. Le
quotient $X/G$ est le schéma torique $\Spec k[x^n,x^{n-1}y, \cdots, xy^{n-1}, y^n]$.

Plus généralement, soient $n$ un entier $\ge 1$, $G$ un groupe abélien d'ordre $n$, $S$ un schéma
noethérien régulier au-dessus de $\Spec \ZZ[1/n,\mu_n]$, $E$ un $\OO_{S}$-module localement
libre de rang fini, muni d'une action linéaire de $G$, $X$ le fibré vectoriel $V(E)$. Pour chaque
caractère $\chi : G \ra \mu_n$, notons $L_{\chi}$ le $G$-$\OO_S$-module correspondant.
L'homomorphisme canonique $G$-équivariant $$ \bigoplus_{\chi} L_{\chi} \otimes E_{\chi} \ra E, $$ où
$E_{\chi} = \mathscr{H}om_G(L_{\chi},E)$ et $\chi$ parcourt le groupe des caractères de $G$, est un
isomorphisme. Il définit une décomposition $G$-équivariante $$ X \iso \oplus_{\chi} X_{\chi}, $$ où
$X_{\chi} = V(E_{\chi})$, muni de l'action de $G$ à travers $\chi$. En particulier, $X^G = X_1$, où
$1 : G \ra \mu_n$ est le caractère trivial. Supposons $S$ local, $S = \Spec A$. Pour chaque $\chi \in
\mathrm{Hom}(G,\mu_n)$, choisissons une base $(t_i)_{i \in I_{\chi}}$ de $\mathrm{Hom}(E_{\chi},A)$,
de sorte que $X_{\chi} = \Spec A[(t_i)_{i \in I_{\chi}}]$, avec $gt_i = \chi(g)t_i$ pour $g \in G$, $i
\in I_{\chi}$. Le couple formé de $X$ et du diviseur à croisements normaux (relatifs à $S$) $Y =
\sum_{\chi, i \in I_{\chi}}Y_i$, où $Y_i = (t_i = 0)$ pour $i \in I_{\chi}$, vérifie les conditions
de \ref{6-dcn stricts} (et les vérifie d'ailleurs fibre à fibre).  

(b) Soient $k$ un corps algébriquement clos d'exposant caractéristique $p$, $n$ un entier $\ge 2$
tel que $(2n,p) = 1$, $G$ le groupe diédral $D_n = ⟨s,r : s^2 = 1, r^n = 1, srs = r^{-1}⟩$. Soit
$\zeta \in k$ une racine primitive $n$-ième de 1. Soit $\rho : G \ra \mathrm{GL}(E)$  la
représentation de degré 2 induite du caractère $\chi$ de $\mu_n \subset G$ tel que $\chi(r) = \zeta$
: $\rho(s) = \begin{pmatrix} 0 & 1 \\ 1 & 0 \\  \end{pmatrix} $, $\rho(r) = \begin{pmatrix} \zeta &
0 \\ 0 & \zeta^{-1} \\ \end{pmatrix}$. Soit $X$ le $G$-schéma $V(E) = \Spec k[u,v]$, $s(u) = v, r(u) =
\zeta u$. Pour $0 \le i \le n-1$, notons $Z_i \subset X$ la droite $v = \zeta^iu$, et $Z = \bigcup_{0
\le i \le n-1}Z_i$. La stratification par l'inertie comporte $n+2$ strates : $X-Z$, où $G$ opère
librement, $Z_i - \{0\}$ ($0 \le i \le n-1$), où le groupe d'inertie est d'ordre 2 (de générateur
$sr^i$), et $\{0\} = X^G$. 

Pour $n = 2$, $G = (\ZZ/2\ZZ)^2$, $Z$ est un diviseur à croisements normaux $G$-strict, et le couple
$(X,Y)$ vérifie les conditions de \ref{6-dcn stricts}. Pour $n > 2$, $Z$ n'est plus un diviseur à
croisements normaux, et l'inertie en $\{0\}$ n'est plus abélienne. Soient $f : X' \ra X$ l'éclaté de
$\{0\}$ dans $X$, $E = f^{-1}(0)$ le diviseur exceptionnel, $Z' = \bigcup_{0 \le i \le n-1}Z'_i$ le
transformé strict de $Z$. Alors $G$ opère de façon naturelle sur $X'$, la projection $f$ est
$G$-équivariante, et $Y' = f^{-1}(Z) = E \cup Z'$ est un diviseur à croisements normaux $G$-strict.
Le couple $(X',Y')$ vérifie les conditions de \ref{6-dcn stricts}. La stratification de $X'$ par
l'inertie se compose des strates $Z'_i$, où l'inertie est un groupe à deux éléments, et de $X'-Z'$,
où $G$ opère librement. La stratification canonique associée à $Y'$ la raffine : $X'^{(0)} = X' -
Y'$, $X'^{(1)} = Y' - \bigcup_{0 \le i \le n-1}(E \cap Z'_i)$, $X'^{(2)} = \bigcup_{0 \le i \le
n-1}(E \cap Z'_i)$.  \subsection{}\label{6-éclatements}\label{6-eclatements} La construction
précédente, qui rend les inerties abéliennes, se généralise. Soit $X$ un schéma noethérien régulier,
séparé, muni d'une action modérée d'un groupe fini $G$, et soit $Y$ un diviseur à croisements
normaux $G$-strict. Si $H$ est un sous-groupe de $G$, $X^H$ est régulier (et séparé), donc il en est
de même de l'éclaté $X' = \Ecl_{X^H}(X)$ de $X$ le long de $X^H$. Le normalisateur $N = N_G(H)$ de $H$
dans $G$ stabilise $X^H$, donc agit sur $X'$, et le morphisme $f : X' \ra X$ est équivariant
relativement à $N \ra G$. De plus, $f^{-1}(X^H)$ est un diviseur régulier dans $X'$. Si $D$ est une
composante de $Y$, comme $D$ est $H$-stable, $D \times_X X^H = D^H$ est régulier, et le transformé
strict $\tilde D = \Ecl_{D^H}(D)$ est un diviseur régulier croisant $f^{-1}(X^H)$ transversalement. Il
s'ensuit que le transformé total réduit $Y' = f^{-1}(Y)_{\mathrm{red}}$ est un diviseur à
croisements normaux $N$-strict dans $X'$.  \begin{proposition}\label{6-inerties abéliennes} Sous les
hypothèses de \ref{6-éclatements}, soit $\overline{x}$ un point géométrique de $X$ en lequel le
groupe d'inertie $G_{\overline{x}}$ n'est pas abélien, et soit $H$ le sous-groupe des commutateurs
$(G_{\overline{x}},G_{\overline{x}})$. Alors $G_{\overline{x}} = N_{G_{\overline{x}}}(H)$ agit sur
$X' = \Ecl_{X^H}(X)$, et en chaque point géométrique $\overline{y}$ de $X'$ au-dessus de
$\overline{x}$, le groupe d'inertie $(G_{\overline{x}})_{\overline{y}}$ est strictement plus petit
que $G_{\overline{x}}$.  \end{proposition} En effet, le point $\overline{y}$ correspond à une droite
$L$ dans la fibre en $\overline{x}$ du fibré normal $T_{\overline{x}}/T_{\overline{x}}^H$ de $X^H$
dans $X$, où $T_{\overline{x}} = T_{\overline{x}}(X)$. Supposons que $\overline{y}$ soit fixe sous
$G_{\overline{x}}$. Alors $G_{\overline{x}}$ agit sur $L$ par un caractère, donc $H$ agit
trivialement sur $L$. Or $(T_{\overline{x}}/T_{\overline{x}}^H)^H = 0$, contradiction. (Noter que
cet argument montre en particulier que, si $H \ne \{1\}$, $X^H$ est de codimension $\ge 2$ dans $X$
en $\overline{x}$.)


\ifx\danslelivre\undefined
\bibliography{../fichiers-configuration-latex/bibliographie}
\bibliographystyle{../fichiers-configuration-latex/smfalpha}
\end{document}
\fi
